\section*{NeurIPS Paper Checklist}

\begin{enumerate}

\item {\bf Claims}
    \item[] Question: Do the main claims made in the abstract and introduction accurately reflect the paper's contributions and scope?
    \item[] Answer: \answerYes{}
    \item[] Justification: The abstract and introduction state the benchmark contribution, the latent-adaptation scope, the heterogeneous scaling pattern, and the interface-reliability finding in a way that matches the reported results and stated caveats; see Abstract, Section 1, and Section 5.

\item {\bf Limitations}
    \item[] Question: Does the paper discuss the limitations of the work performed by the authors?
    \item[] Answer: \answerYes{}
    \item[] Justification: The main paper contains a dedicated Limitations section covering provider heterogeneity, interface confounding, adaptation-metric sharpness, and synthetic-world scope; see Section 6.

\item {\bf Theory assumptions and proofs}
    \item[] Question: For each theoretical result, does the paper provide the full set of assumptions and a complete (and correct) proof?
    \item[] Answer: \answerNA{}
    \item[] Justification: The paper is an empirical benchmark paper and does not present formal theorems or proofs. The equations define benchmark mechanics and grading, not new theoretical results.

\item {\bf Experimental result reproducibility}
    \item[] Question: Does the paper fully disclose all the information needed to reproduce the main experimental results of the paper to the extent that it affects the main claims and/or conclusions of the paper (regardless of whether the code and data are provided or not)?
    \item[] Answer: \answerYes{}
    \item[] Justification: The manuscript and release bundle specify the frozen protocol, task scope, seed counts, split structure, model scope, regeneration commands, resolved manifest, and artifact layout; exploratory hosted API probes are labeled separately and their raw logs are included under \texttt{outputs/api-probe-2026-05-07}. See Sections 2--4 and the release files \texttt{paper/release/README.md} and \texttt{paper/release/REPRODUCTION.md}.

\item {\bf Open access to data and code}
    \item[] Question: Does the paper provide open access to the data and code, with sufficient instructions to faithfully reproduce the main experimental results, as described in supplemental material?
    \item[] Answer: \answerYes{}
    \item[] Justification: The submission now includes an anonymous reviewer archive that packages the runnable benchmark code, the frozen \texttt{2026-04-18} analysis root, the exact raw run directories referenced by \texttt{metadata/run\_manifest.json}, the exploratory hosted API probe logs, and reproduction instructions; see Section 4 and the release files \texttt{paper/release/README.md}, \texttt{paper/release/REPRODUCTION.md}, and \texttt{paper/release/REVIEW\_BUNDLE.md}.

\item {\bf Experimental setting/details}
    \item[] Question: Does the paper specify all the training and test details (e.g., data splits, hyperparameters, how they were chosen, type of optimizer) necessary to understand the results?
    \item[] Answer: \answerYes{}
    \item[] Justification: The paper and release package specify the benchmark tasks, hidden-shift configuration, split structure, seed counts, scoring definitions, serving path, and model roles needed to interpret the evaluation; hosted API probes additionally disclose provider route, model identifiers, split coverage, and their exploratory status. See Sections 2--4 and the release protocol files.

\item {\bf Experiment statistical significance}
    \item[] Question: Does the paper report error bars suitably and correctly defined or other appropriate information about the statistical significance of the experiments?
    \item[] Answer: \answerYes{}
    \item[] Justification: The main result and baseline tables report 95\% bootstrap confidence intervals on the headline aggregates, and the text explains the role of native-only estimates and effective sample size; see Section 4 and the generated result tables.

\item {\bf Experiments compute resources}
    \item[] Question: For each experiment, does the paper provide sufficient information on the computer resources (type of compute workers, memory, time of execution) needed to reproduce the experiments?
    \item[] Answer: \answerNo{}
    \item[] Justification: The release package reports provider path, token counts, preserved wall-clock timing, and a validation-host note in \texttt{paper/release/COMPUTE.md}, but the exact original worker inventory for the full frozen runs was not preserved in the manifest. We therefore keep the conservative No answer rather than over-claim hardware-normalized reproducibility.

\item {\bf Code of ethics}
    \item[] Question: Does the research conducted in the paper conform, in every respect, with the NeurIPS Code of Ethics \url{https://neurips.cc/public/EthicsGuidelines}?
    \item[] Answer: \answerYes{}
    \item[] Justification: To the best of our knowledge, yes. The benchmark is synthetic, does not involve human subjects, and the paper and datasheet explicitly restrict inappropriate deployment-style uses; see Section 6 and \texttt{paper/release/DATASHEET.md}.

\item {\bf Broader impacts}
    \item[] Question: Does the paper discuss both potential positive societal impacts and negative societal impacts of the work performed?
    \item[] Answer: \answerYes{}
    \item[] Justification: The release datasheet and the paper's limitations discussion address appropriate uses, inappropriate uses, and the risk of over-interpreting a synthetic benchmark as a real deployment simulator; see Section 6 and \texttt{paper/release/DATASHEET.md}.

\item {\bf Safeguards}
    \item[] Question: Does the paper describe safeguards that have been put in place for responsible release of data or models that have a high risk for misuse (e.g., pre-trained language models, image generators, or scraped datasets)?
    \item[] Answer: \answerNA{}
    \item[] Justification: The paper releases a synthetic benchmark and frozen evaluation artifacts, not a high-risk generative model or scraped dataset. The benchmark also includes intended-use and limitation guidance in the release datasheet.

\item {\bf Licenses for existing assets}
    \item[] Question: Are the creators or original owners of assets (e.g., code, data, models), used in the paper, properly credited and are the license and terms of use explicitly mentioned and properly respected?
    \item[] Answer: \answerYes{}
    \item[] Justification: The release package now includes a consolidated asset-and-terms inventory for the evaluated external model families and serving dependencies, with official source links and notes on local alias handling; see \texttt{paper/release/THIRD\_PARTY\_ASSETS.md}.

\item {\bf New assets}
    \item[] Question: Are new assets introduced in the paper well documented and is the documentation provided alongside the assets?
    \item[] Answer: \answerYes{}
    \item[] Justification: The benchmark release includes documented code, a frozen artifact bundle, reproduction instructions, a datasheet, and Croissant metadata; see Section 4 and the files under \texttt{paper/release/}.

\item {\bf Crowdsourcing and research with human subjects}
    \item[] Question: For crowdsourcing experiments and research with human subjects, does the paper include the full text of instructions given to participants and screenshots, if applicable, as well as details about compensation (if any)?
    \item[] Answer: \answerNA{}
    \item[] Justification: The paper does not involve crowdsourcing or human-subject experiments.

\item {\bf Institutional review board (IRB) approvals or equivalent for research with human subjects}
    \item[] Question: Does the paper describe potential risks incurred by study participants, whether such risks were disclosed to the subjects, and whether Institutional Review Board (IRB) approvals (or an equivalent approval/review based on the requirements of your country or institution) were obtained?
    \item[] Answer: \answerNA{}
    \item[] Justification: The paper does not involve crowdsourcing or human-subject experiments, so IRB approval is not applicable.

\item {\bf Declaration of LLM usage}
    \item[] Question: Does the paper describe the usage of LLMs if it is an important, original, or non-standard component of the core methods in this research? Note that if the LLM is used only for writing, editing, or formatting purposes and does \emph{not} impact the core methodology, scientific rigor, or originality of the research, declaration is not required.
    \item[] Answer: \answerYes{}
    \item[] Justification: LLMs are the benchmarked systems in the experimental evaluation, including locally served models and hosted OpenRouter probes. The benchmark construction, deterministic scoring, and analysis pipeline do not rely on an auxiliary LLM-based labeling or generation procedure.

\end{enumerate}
